% 术语表 Glossary
\chapter*{术语表}
\addcontentsline{toc}{chapter}{术语表}
\chaptermark{术语表}

\noindent\textit{Glossary}

以下条目按英文字母顺序排列。英文术语保留原文；定义已译为中文，命令名与技术标识符保持英文。

\begin{description}[style=nextline, leftmargin=3.2cm, labelwidth=2.8cm, font=\normalfont]

\item[标注（\textit{annotation}）]
附着于设计中对象的信息片段，例如附着于网上的电容值；或将此类信息附着到设计中对象的过程。

\item[反标（\textit{back-annotate}）]
使用提取及其他后处理信息更新电路设计的过程，这些信息反映设计依赖于实现的特性，例如引脚选择、组件位置或寄生电气特性。
反标允许对最终电路做更准确的时序分析。数据由布局后的另一工具生成并传递到综合环境。
例如，设计数据库可用实际互连延时更新；这些延时在布局布线之后——确切互连长度已知之后——计算。

\item[单元（\textit{cell}）]
见实例（instance）。

\item[Clock（\textit{clock}）]
具有周期行为的定时脉冲源。clock 通过控制数字电路中的 sequential element（如 flip-flop 与 register）来同步数据信号传播。用 \cmd{create\_clock} 定义 clock。

用 \cmd{create\_clock} 创建的 clock 默认忽略 clock network 的 delay effect。因此，为进行准确的 timing analysis，需要用 clock latency 与 clock skew 描述 clock network。另见 clock latency 与 clock skew。

\item[时钟门控（\textit{clock gating}）]
用逻辑（非反相器或缓冲器）控制时钟信号，以在选定时刻关断时钟信号或修改时钟脉冲特性。

\item[Clock latency（\textit{clock latency}）]
clock signal 从 clock source 传播到设计中特定点所需的时间。clock latency 是 source latency 与 network latency 之和。

source latency 是从实际 clock origin 到设计中 clock definition point 的传播时间。network latency 是从 clock definition point 到第一个 register clock pin 的传播时间。

用 \cmd{set\_clock\_latency} 命令指定时钟延迟。

\item[Clock skew（\textit{clock skew}）]
同一 clock domain 内或不同 clock domain 之间，各 register 处 clock arrival time 的最大差值。本指南在此处用 clock uncertainty 表征该 skew 特性，并使用 \cmd{set\_clock\_uncertainty} 为一个或多个 clock network 指定该值；clock skew 是实际到达时间之差，不能一般性地与包含 jitter、建模裕量等成分的 clock uncertainty 等同。

\item[时钟源（\textit{clock source}）]
将时钟波形施加到设计的引脚或端口。时钟信号到达其所有源传递扇出中的寄存器。时钟可有多个源。

对 \cmd{create\_clock} 使用 \texttt{source\_object} 选项指定时钟源。

\item[时钟树（\textit{clock tree}）]
时钟源与该源传递扇出中寄存器之间的组合逻辑。时钟树亦称时钟网络，由厂商根据一个时钟域内或时钟域之间寄存器处的物理布局数据综合。

\item[Clock uncertainty（\textit{clock uncertainty}）]
本指南此处将其用于表征 clock skew；见 clock skew。实际 SDC 分析中，clock uncertainty 是施加于 clock relationship 的 timing margin，不应与实际 clock arrival time 之差混为一谈。

\item[核（\textit{core}）]
用作 ASIC 设计构建块的预设计逻辑块。

\item[Critical path（\textit{critical path}）]
电路中 delay 最长的 path。电路速度取决于最慢的 register-to-register delay。clock period 不能短于该 delay，否则 signal 无法及时到达下一级 register 并被捕获。

\item[当前设计（\textit{current design}）]
活动设计（正在处理的设计）。多数命令特定于当前设计，即在当前设计上下文中操作。用 \cmd{current\_design} 命令指定当前设计。

\item[当前实例（\textit{current instance}）]
设计层次中实例专用命令默认操作的实例。用 \cmd{current\_instance} 命令指定当前实例。

\item[Datapath（\textit{datapath}）]
用加法器、乘法器、移位器与比较器等算术运算符操纵数据信号的逻辑电路。

\item[设计约束（\textit{design constraints}）]
设计者对设计性能目标的规格，即综合在其下执行的时序与环境限制。DC Explorer 用这些约束——例如低功耗、小面积、高速度或最小成本——指导设计优化以满足面积与时序目标。

设计约束分两类：设计规则约束与设计优化约束。
\begin{itemize}
  \item 设计规则约束在工艺库中提供。为保证所制造电路正常工作，不得违例。
  \item 设计优化约束定义时序与面积优化目标。
\end{itemize}
DC Explorer 按两组约束优化设计综合；但设计规则约束优先级更高。

\item[伪路径（\textit{false path}）]
您不希望 DC Explorer 在时序分析期间考虑的路径。例如两个永不同时使能的多路复用块之间的路径，即不能传播信号的路径。

用 \cmd{set\_false\_path} 命令按路径禁用基于时序的综合。该命令移除指定路径上的时序约束。

\item[扇入（\textit{fanin}）]
驱动端点引脚、端口或网（亦称 sink）的引脚。若存在经组合逻辑从该引脚到 sink 的时序路径，则该引脚在 sink 的扇入中。扇入追踪从寄存器时钟引脚或有效起点开始。扇入亦称传递扇入（transitive fanin）。

用 \cmd{report\_transitive\_fanin} 命令报告指定 sink 引脚、端口或网的扇入。

\item[扇出（\textit{fanout}）]
由源引脚、端口或网驱动的引脚。若存在经组合逻辑从源到该引脚或端口的时序路径，则该引脚在源的扇出中。扇出追踪在寄存器数据引脚或有效端点处停止。扇出亦称传递扇出（transitive fanout）或时序扇出。

用 \cmd{report\_transitive\_fanout} 命令报告指定源引脚、端口或网的扇出。

\item[扇出负载（\textit{fanout load}）]
表示对总扇出数值贡献的无量纲数。扇出负载不同于负载（电容值）。

DC Explorer 通过将 \texttt{fanout\_load} 属性与每个输入引脚关联、将 \texttt{max\_fanout} 属性与单元每个输出（驱动）引脚关联来建模扇出限制，并确保扇出负载之和小于 \texttt{max\_fanout} 值。

\item[展平（\textit{flatten}）]
将设计的组合逻辑路径转换为两级积之和表示。展平期间，DC Explorer 从设计中移除所有中间项及因此所有相关逻辑结构。展平基于约束。

\item[前向标注（\textit{forward-annotate}）]
将数据从综合环境传输到设计流程后续所用其他工具。例如，可将标准延时格式（SDF）中的延时与约束数据从综合环境传输以指导布局布线工具。

\item[生成时钟（\textit{generated clock}）]
由集成电路内部生成的时钟信号；非直接来自外部源。例如由系统时钟生成的二分频时钟。用 \cmd{create\_generated\_clock} 命令定义生成时钟。

\item[Hold time（\textit{hold time}）]
clock active edge 之后，data pin 上的 signal 必须保持稳定的时间。hold time 为到达 cell data pin 的 path 建立 minimum-delay requirement。

本指南给出的 hold 检查关系式为：
\[
\text{hold} = \text{max clock delay} - \text{min data delay}
\]

\item[理想时钟（\textit{ideal clock}）]
被视为在时钟网络中传播时无延时的时钟。理想时钟类型是 DC Explorer 的默认。可用 \cmd{set\_clock\_latency} 与 \cmd{set\_propagated\_clock} 命令覆盖默认行为，以获得非零时钟网络延时并指定关于时钟网络延时的信息。

\item[理想网（\textit{ideal net}）]
被赋予理想时序条件的网——即延迟、转换时间与电容赋值为零。此类网免于时序更新、延时优化与设计规则修复。将打算单独综合的某些高扇出网（如 scan-enable 与复位网）定义为理想网可缩短运行时间。

用 \cmd{set\_ideal\_net} 命令将网指定为理想网。

\item[Input delay（\textit{input delay}）]
指定从 clock edge 到 signal 到达指定 input port 的 minimum 或 maximum delay 的 constraint。

用 \cmd{set\_input\_delay} 命令相对于指定时钟信号在引脚或输入端口上设置输入延时。

\item[实例（\textit{instance}）]
内存中已加载引用（库组件或设计）在电路中的出现；每个实例有唯一名称。设计可包含多个实例；每个实例指向同一引用，但有唯一名称以与其他实例区分。实例亦称单元（cell）。

\item[叶单元（\textit{leaf cell}）]
逻辑设计的基本单元。叶单元不能分解为更小的逻辑单元。例如 NAND 门与反相器。

\item[链接库（\textit{link library}）]
DC Explorer 用于解析单元引用的工艺库。链接库可包含工艺库与设计文件。链接库还包含已映射网表中单元（库单元以及子设计）的描述。

链接库包括本地链接库（\texttt{local\_link\_library} 属性）与系统链接库（\varname{link\_library} 变量）。

\item[多周期路径（\textit{multicycle path}）]
数据从起点传播到端点需要超过一个时钟周期的路径。

用 \cmd{set\_multicycle\_path} 命令指定 DC Explorer 应用于确定数据在特定端点何时需要的时钟周期数。

\item[网表（\textit{netlist}）]
以 ASCII 或二进制格式描述电路原理图的文件——网表包含设计中电路元件与互连的列表。网表传输是在设计系统或工具之间移动设计信息的最常用方式。

\item[工作条件（\textit{operating conditions}）]
设计遇到的工艺、电压与温度范围。DC Explorer 按工艺、电压与温度曲线上的工作点优化设计，并根据工作条件缩放单元与线延时。

默认情况下，工作条件在工艺库的 \texttt{operating\_conditions} 组中指定。

\item[优化（\textit{optimization}）]
逻辑综合过程中 DC Explorer 尝试实现最能满足设计功能、时序与面积要求的工艺库单元组合的步骤。

\item[Output delay（\textit{output delay}）]
指定从 output port 到捕获该 output 数据的 sequential element 的 minimum 或 maximum delay 的 constraint。该 constraint 确定 signal 必须在 output port 可用的时间，以满足 sequential element 的 setup/hold requirement。

用 \cmd{set\_output\_delay} 命令相对于指定时钟信号在引脚或输出端口上设置输出延时。

\item[Pad 单元（\textit{pad cell}）]
位于芯片边界、允许与芯片外集成电路连接或通信的特殊单元。

\item[路径组（\textit{path group}）]
相关路径的组，由 \cmd{create\_clock} 命令隐式分组或由 \cmd{group\_path} 命令显式分组。默认情况下，端点由同一时钟打入的路径分配到同一路径组。

\item[引脚（\textit{pin}）]
单元上提供输入与输出连接的部分。引脚可为双向。子设计的端口在父设计中是引脚。

\item[传播时钟（\textit{propagated clock}）]
经时钟网络产生延时的时钟。传播时钟用于确定寄存器时钟引脚处的时钟延迟。由传播时钟打入的寄存器，其边沿时间因从时钟源到寄存器时钟引脚的路径延时而偏斜。

用 \cmd{set\_propagated\_clock} 命令指定时钟延迟经时钟网络传播。

\item[实时钟（\textit{real clock}）]
具有源的时钟，意味着其波形施加到设计中的引脚或端口。通过使用 \cmd{create\_clock} 命令并包含端口或引脚的源列表来创建实时钟。实时钟可以是理想的或传播的。

\item[引用（\textit{reference}）]
可用作构建更大电路的元素的库组件或设计。引用的结构可以是简单逻辑门或更复杂的设计（RAM 核或 CPU）。设计可包含引用的多次出现；每次出现是一个实例。亦见实例。

\item[RTL（\textit{RTL}）]
RTL（register transfer level，寄存器传输级）是数字电子电路的寄存器级描述。在数字电路中，寄存器在时钟周期之间存储中间信息；因此 RTL 描述所存储的中间信息、其在设计中的存储位置，以及如何通过设计传输。RTL 在一组寄存器之间的数据流级别建模电路行为。此抽象级别通常几乎不含时序信息，除对一组时钟边沿与特性的引用外。

\item[Setup time（\textit{setup time}）]
clock active edge 之前，data pin 上的 signal 必须保持稳定的时间。setup time 为到达 cell data pin 的 path 建立 maximum-delay requirement。

本指南给出的 setup 检查关系式为：
\[
\text{setup} = \text{max data delay} - \text{min clock delay}
\]

\item[Slack（\textit{slack}）]
表示 mapped design 中 path endpoint 处 data arrival time 相对于 required arrival time 的 timing margin。slack 可为正、负或零；具体差值方向取决于所执行的 setup（max-delay）或 hold（min-delay）检查，但符号含义一致。

正 slack 表示仍满足 timing constraint；对 setup path，它表示 path delay 在不违例约束的情况下可增加的量。负 slack 表示 timing violation；对 setup path，其绝对值表示为满足约束必须减少的 path delay。hold slack 应按 min-delay 检查解释，不能直接套用 setup path 的“减少 delay”表述。

\item[结构化（\textit{structuring}）]
向设计添加中间变量与逻辑结构，可导致减小设计面积。结构化基于约束。最适合应用于非关键时序路径。

默认情况下，DC Explorer 对设计做 structuring。

\item[综合（\textit{synthesis}）]
从输入 IC 设计生成基于工艺库的优化门级网表的软件过程。综合包括读取 HDL 源代码并基于该描述优化设计。

\item[符号库（\textit{symbol library}）]
包含特定 ASIC 库中所有单元原理图符号的库。DC Explorer 用符号库生成设计原理图。可用 Design Vision 查看设计原理图。

\item[目标库（\textit{target library}）]
DC Explorer 在优化期间映射到的工艺库。目标库包含用于生成网表的单元以及设计系统工作条件的定义。

\item[工艺库（\textit{technology library}）]
综合过程中 DC Explorer 可用的 ASIC 单元库。工艺库可包含每个 ASIC 单元的面积、时序、功耗与功能信息。每个库的工艺特定于特定 ASIC 厂商。

\item[时序例外（\textit{timing exception}）]
对 DC Explorer 假定的默认（单周期）时序行为的例外。为使 DC Explorer 正确分析电路，必须指定设计中不符合默认行为的每条时序路径。时序例外的例子包括伪路径、多周期路径，以及需要与默认计算时间不同的特定最小或最大延时时间的路径。

\item[时序路径（\textit{timing path}）]
规定数据通过设计传播的点到点序列。数据在起点由时钟边沿发射，经组合逻辑元件传播，并在端点由另一时钟边沿捕获。时序路径的起点是输入端口或时序元件的时钟引脚。时序路径的端点是输出端口或时序元件的数据引脚。

\item[转换延时（\textit{transition delay}）]
由驱动引脚改变电压状态所需时间引起的时序延时。

\item[取消分组（\textit{ungroup}）]
移除设计中的层次级别。Ungrouping 将给定层次级别的子设计合并到父单元或设计中。

用 \cmd{ungroup} 命令，或对 \cmd{compile} 使用 \texttt{auto\_ungroup} 选项来 ungroup 设计。

\item[唯一化（\textit{uniquify}）]
解析内存中对同一设计的多个单元引用。

唯一化过程为每个引用原始设计的已实例化单元创建具有唯一设计名的唯一设计副本。

\item[虚时钟（\textit{virtual clock}）]
存在于系统中但不属于该块的时钟。虚时钟不打入当前设计内任何时序器件，也不与引脚或端口关联。用虚时钟作为相对于块外时钟指定输入与输出延时的参考。

用不带关联引脚或端口列表的 \cmd{create\_clock} 命令创建虚时钟。

\end{description}
